\section{The BX3 Three-Layer Model}\label{sec:three-layer-model}

The \bxthree{} Framework's three-layer model is the architectural substrate
within which all five pillars operate as mutually reinforcing structural
necessities.  The model organizes every node into three immutable functional
layers --- \purpose{}, \bounds{}, and \fact{} --- with strictly defined roles,
interaction constraints, and enforcement primitives.

This section defines each layer, establishes the functional separation that
makes the layer roles distinct, and defines the three invariants that govern
inter-layer interaction.

% --------------------------------------------------------------
\subsection{Layer Definitions}\label{sec:layer-definitions}
% --------------------------------------------------------------

The \purpose{} layer carries intent, judgment, and final authority.  It is the
layer at which the human principal's objectives are translated into operational
directives, and at which the operating range $R(n)$ for each node $n$ is
defined.  It is the exclusive terminus for any exception state that exceeds
machine-actor provisioned authority.  The layer is non-delegable: no machine
actor possesses the judgment capacity to serve as a substitute terminus for
exceptions that exceed its own operating range.  This non-delegability is not a
policy constraint; it is a structural property of the layer's function.  The
\purpose{} layer carries the principal's intent, and no agent operating under
delegated authority can possess the principal's intent in sufficient
completeness to serve as its substitute terminus.

The \bounds{} layer carries bounded reasoning and constrained execution within
$R(n)$.  It evaluates whether proposed actions fall within the node's delegated
authority by reference to the authorization ledger maintained by the \fact{}
layer.  When it encounters a condition outside $R(n)$ that it cannot resolve
within its own scope, it triggers the Bailout Protocol compulsorily --- the
trigger is a structural consequence of the bounds check, not a discretionary
decision by the agent.  The \bounds{} layer cannot choose to absorb the
exception; the architectural separation between \purpose{} and \bounds{}
ensures that any condition outside $R(n)$ is, by definition, beyond the
\bounds{} layer's adjudication authority.

The \fact{} layer carries deterministic enforcement and forensic auditability.
It maintains the append-only Forensic Ledger recording every protocol event
with full attribution.  It enforces the transition relations of all pillar
state machines and has exclusive write authority over the Forensic Ledger.  No
machine actor may modify, delete, or suppress entries in the Forensic Ledger;
the layer's function is to enforce the record rather than to interpret it.  It
is the only layer that may issue a definitive halt to autonomous action, and it
does so by refusing to approve any proposed action that lacks a valid
authorization entry in the ledger.

% --------------------------------------------------------------
\subsection{Functional Separation Invariants}\label{sec:layer-invariants}
% --------------------------------------------------------------

The three layers are not merely a conceptual taxonomy; they are enforced
architectural divisions with distinct execution semantics captured in three
structural invariants.

\begin{invariant}[Purpose-Bounds Separation]\label{inv:pb-sep}
No machine actor operating in the \bounds{} layer may issue binding resolution
directives that close a protocol run, define or modify the operating range
$R(n)$, or commit entries to the authorization ledger bearing a directive from
the \purpose{} layer.  The \bounds{} layer executes bounded reasoning within
$R(n)$; it does not carry intent.
\end{invariant}

\begin{invariant}[Bounds-Fact Separation]\label{inv:bf-sep}
No node in the \bounds{} layer may modify the transition relations of any
pillar state machine, alter the authorization ledger, or suppress entries in
the Forensic Ledger.  The \fact{} layer enforces these relations
deterministically; the enforcement is independent of the \bounds{} layer's
operational state.
\end{invariant}

\begin{invariant}[Fact-Purpose Convergence]\label{inv:fp-converge}
The \fact{} layer enforces the exclusive authority of the \purpose{} layer
over binding resolution directives by rejecting any directive whose issuer is
not verified as $h(n)$.  The \fact{} layer does not carry intent; it carries
the enforcement mechanism by which the \purpose{} layer's authority is made
structurally operative.
\end{invariant}

INV\textsubscript{pb-sep} and INV\textsubscript{bf-sep} together establish the
layer separation that makes the \purpose{} layer the sole terminus for
unresolved exceptions: the \bounds{} layer can detect and attempt to resolve
exceptions but cannot close the protocol; the \fact{} layer can enforce and
record but cannot adjudicate.  INV\textsubscript{fp-converge} closes the loop
by ensuring that the \fact{} layer's enforcement is定向 toward the
\purpose{} layer's authority rather than toward any machine actor's claim.

Together, the three invariants constitute a \emph{separation of mandate}
property: intent (Purpose), execution (Bounds), and enforcement (Fact) are
architecturally distinct, and no layer may exercise the mandate of another.
This separation is what makes the \bxthree{} Framework's accountability
guarantees structurally operative rather than aspirational.

% --------------------------------------------------------------
\subsection{The Accountability Chain}\label{sec:accountability-chain}
% --------------------------------------------------------------

The three-layer model establishes a two-stage accountability chain for every
\bxthree{} node $n$:

\begin{enumerate}\itemsep=0pt
\item[(i)] The \purpose{} layer of $n$ defines $R(n)$ and provisions the
\bounds{} layer with the authorization scope within which $n$ may act
autonomously.  This provisioning is recorded as an authorization ledger entry
by the \fact{} layer.

\item[(ii)] If the \bounds{} layer of $n$ encounters a condition outside $R(n)$
that it cannot resolve, the Bailout Protocol compels propagation of the
exception to $h(n)$ via the parent chain.  The \purpose{} layer of $h(n)$
receives the exception and issues a binding directive, which is recorded in
the Forensic Ledger by the \fact{} layer of the originating node.
\end{enumerate}

This two-stage chain is the structural mechanism by which the upstream
accountability guarantee is realized: every consequential decision in any
\bxthree{} node is either (a) authorized by the \purpose{} layer and enforced
by the \fact{} layer, or (b) escalated to $h(n)$ through the Bailout
Protocol when $R(n)$ is exceeded.  No decision proceeds from (b) without a
binding directive from $h(n)$ recorded in the Forensic Ledger.  The chain
is append-only and the parent pointer $\alpha(n)$ is immutable, ensuring
that the chain can be reconstructed at any time for any node.